% =====================================================================
%  R. Nielsen, EFTER DØDEN, BIDRAG TIL EN BELYSNING AF SELVHEDSLOVEN
%  In: For Ide og Virkelighed. Et Tidsskrift, udgivet af R. Nielsen, Bj. Bjørnson og Rud. Schmidt.  Aargang 1871. Andet Bind (Bind 5).
%  Kjøbenhavn: Otto B. Wroblewskys Forlag; G. S. Wibes Bogtrykkeri.  Bind 5, pp. 193-223.  Antiqua.
%
%  TRANSCRIPTION -- GENERATED; DO NOT EDIT BY HAND.  Rasmus Nielsen's eight essays in the volumes 1870-1873 are built by ONE
%  pipeline (ebook-method/for-idee-1870-73/) into a combined working file, which split7.py cuts into seven transcriptions
%  (Aandsvirkelighed and its sequel, Videre om Aandsvirkelighed, are one: Hans, 2026-10-13).
%  Rebuild from the repo root:   sh ebook-method/for-idee-1870-73/rebuild.sh
%
%  STATUS: COMPLETE (2026-10-13): the volume's title page and the whole essay, Bind 5, pp. 193-223, with its notes.
%
%  EDITION.  The only printing.  Text of record: Nasjonalbiblioteket's scan (nb.no) of the volume, ~/bibliotek/Nielsen, Rasmus/1871-2-for-idee-og-virkelighed.pdf (NB item 9bae6062efbd4219d60f56932f2a37a9) (SCANS.tsv),
%  page images fetched by nbfetch.py, 300 dpi colour, NO text layer; printed page N = PDF page N + 4.  No e-book, no
%  Wikisource text, no reprint, no earlier transcription found (checked 2026-10-13: the repo, the web, the Internet Archive,
%  da.wikisource; HathiTrust and the Google Books API could not be queried that day; Google Books holds at least the 1873
%  volume, nielsen/darwinismen's source -- a second copy, not used).
%
%  ACCURACY (2026-10-13; one collation for all eight texts, 231 pp., ~58,000 words).  The DRAFT is NB's own OCR (ALTO, ABBYY
%  FineReader 12), laid over the page images (combine.py) and turned into LaTeX by draft.py; it was collated word by word
%  with an independent tesseract reading of the same images (collate/check.py, witness 'ocr'): 1556 disagreements.
%  check.py's own rules settled 78.  Candidate rules were first scored BY SCRIPT on a text already read at the image
%  (calibrate.py: Et Synspunkt for Darwinismen, same volume as Verdensanskuelse), then measured on a BLIND PILOT of 152 crops
%  (reader A, 52/52 planted controls right); 17 rules passed (NB's glyph faults: o for ø, D/O/U for Ø, i for l, li for h,
%  l for I, opening » read «, words run together or broken apart; tesseract's: aa read an/ag/å, added accents; the
%  printer's space before a mark, house rule) and settled 772; the rest were read at the image BLIND (562 crops on 94
%  sheets, four readers, 188/188 planted controls right), ~160 decided in the main session from the readers' lines and ~40
%  at 300-dpi zooms.  STRUCTURE read at the image (reader K): every footnote whole (30; collated where they stand, then moved
%  to their marks, then compared word by word with the reading: notecmp.py, 9 differences settled at the image), the
%  headings, the section numerals, the essays' last lines and signatures, the three Greek words, the volume title pages.
%  LETTERSPACING (the print's emphasis, \emph): candidate lines by letter gaps and box width, calibrated on the words the
%  word readers had marked spaced (spaced.py): 361 lines read blind (61/62 roman controls), 25 doubtful extents read at a
%  zoom.  SWEEPS after OPEN 0: words run together that both OCRs share (lower-case + capital), the unknown-word sweep (678
%  words, ~20 suspects, 14 read at the image: 4 OCR faults, 7 misprints kept), --moved, the dropped-word list, the
%  W-shortens list of the blind decisions.
%  SAMPLED PAGE CHECK (pagecheck.py, seed 202610137): one page of each essay or part, pp. 2289, 2400, 3305, 4308, 5033,
%  5208, 6514, 8011 of the working numbering (Bind x 1000 + page), read whole and blind, 1965 words: 0 word errors,
%  1 punctuation (a compound's hyphen lost at a page turn: „Ny-|Holland“) and 1 letterspacing („to be or not to be“)
%  error, both mended.  95 % upper bound for word errors: 1.5 per 1000 words.
%  What this cannot see: a misprint both OCRs normalise into a real word.
%  The ø rule (Hans 2026-10-08): o-forms that are no word are set with ø; „stode“ (= stood) is a word and stays.
%  PRINTED AS IS: nothing recorded.
%  Not transcribed: running heads, folios, sheet signatures („For Ide og Virkelighed. 1870. I. 20“), library stamps; the
%  large initial of each essay is set as a plain letter.  Quotation marks as printed (» « in the text, sometimes a closing
%  »; „ “ in some notes).  The notes are numbered ¹), ²) per page, as printed.  Greek as printed (textalpha).
%
%  TOOLS (ebook-method/for-idee-1870-73/): nbalto.py, combine.py, alto.py, ocr.py, draft.py, rules.py, calibrate.py,
%  sheets.py + reads.py + READER-BRIEF.md, spaced.py + spaced.tsv, structsheets.py, notecmp.py, unknown.py, zoom.py,
%  pagecheck.py + PAGECHECK-BRIEF.md, patches.py (PATCHES, PRENOTES, NOTES, POST, MISPRINTS), decisions.tsv, split7.py,
%  rebuild.sh.
% =====================================================================
\documentclass[12pt,a4paper,openany]{book}
\usepackage[utf8]{inputenc}
\usepackage[T1]{fontenc}
\usepackage{libertinus}
\usepackage{libertinust1math}
\usepackage{textalpha}
\usepackage[danish]{babel}
\usepackage[protrusion=true,expansion=true]{microtype}
\usepackage{setspace}
\usepackage[margin=1.3in]{geometry}
\usepackage{fancyhdr}
\usepackage{hyperref}
\hypersetup{pdftitle={Efter Døden}, pdfauthor={Rasmus Nielsen}, hidelinks}
\pagestyle{fancy}
\fancyhf{}
\fancyhead[LE,RO]{\thepage}
\fancyhead[C]{\textit{R. Nielsen: Efter Døden (1871)}}
\setstretch{1.3}
\usepackage{marginnote}
\renewcommand{\thefootnote}{\arabic{footnote})}   % the print numbers each page's notes ¹), ²)
\newcommand{\opage}[1]{\marginnote{\footnotesize\textit{[#1]}}}
% the essay's heading: title (and subtitle), centred, a short rule; the print's large initial is not reproduced
\newcommand{\essayhead}[2]{\begin{center}{\large #1}\ifx&#2&\else\\[0.3em]{\normalsize #2}\fi\\[0.5em]\rule{3em}{0.4pt}\end{center}\medskip\noindent}
\newcommand{\signature}[1]{\par\smallskip\hfill #1\hspace*{2em}\par}

\begin{document}

% --- front matter: the title page of Bind 5 (PDF 1 of 1871-2-for-idee-og-virkelighed.pdf), read at the image (reader K, m1);
% short wavy rules are set as plain rules; the library stamps are not transcribed.
\begin{titlepage}
\begin{center}
\vspace*{3em}
{\Huge For Ide og Virkelighed.}\\[1em]
\rule{8em}{0.4pt}\\[1em]
{\large Et Tidsskrift,}\\[0.8em]
{\small udgivet}\\[0.4em]
{\footnotesize af}\\[0.6em]
{\large R. Nielsen,}\\[0.4em]
{\large Bj. Bjørnson og Rud. Schmidt.}\\[1em]
\rule{8em}{0.4pt}\\[0.6em]
{\large Aargang 1871. Andet Bind.}\\[0.6em]
\rule{8em}{0.4pt}\\[4em]
Kjøbenhavn.\\[0.4em]
Otto B. Wroblewskys Forlag.\\[1.5em]
{\small G. S. Wibes Bogtrykkeri.}
\end{center}
\end{titlepage}

% ===== TEXT: pp. 193-223 =====
\essayhead{Efter Døden,}{Bidrag til en Belysning af Selvhedsloven.}
% --- p. 193 ---
\setcounter{footnote}{0}%
\opage{193}Hvor Døden dog er en underlig Gjæst! Den færdes i
Vrimlen, men passer sit Eget; den gaaer ud og ind i
Menneskenes Huse lig en gammel Betjent, en Slags officiel
Person, med hvis Bestilling Slægten fra Arilds Tid
er bleven fortrolig. Taus og stille fuldbyrder Manden med
Leen sin vante Gjerning, en almindelig Hverdagsforretning,
en Gjerning, som Alle kjende, og som dog, naar
der spørges om Sagen, Ingen kjender, thi Gjerningen selv
er en Gaade, en uigjennemtrængelig Hemmelighed. Den,
der har prøvet, hvad Sorg er, kan tale med om Sorger,
og Sorgernes Erindring vil gjøre ham veltalende; den,
hvem Ulykken engang har ramt, kan tale med om Ulykke,
og Gjenskin af Rædsel og Angst endnu give Glimt i hans
Tale; men den, hvem Døden har ramt, kan ikke tale
med, kan ikke med et Ord forraade, hvad der er skeet:
den Døde tier med hvad han veed; Knokkelmanden har
som for en Forsigtigheds Skyld trykt et Taushedssegl paa
hans Læber.

Ligesom Dødens Gjerning er en Hemmelighed, saaledes
er den Time, det Øieblik, da den vil indfinde sig,
% --- p. 194 ---
\setcounter{footnote}{0}%
\opage{194}i det Hele jo ogsaa en Hemmelighed. Den, der paa
Embeds Vegne reiser omkring i Landet og seer Kasserne
efter, ønsker helst at komme bag paa Vedkommende for
at erfare, om Regnskaberne ere i Orden. Man skulde
næsten troe, at Døden gik i et lignende Ærinde, thi den
holder af at overraske Menneskene. Undertiden, gjennem
Sygdom og Smerter, skeer der Anmeldelse, at nu er Øieblikket
nær, nu kommer den Tause med Leen; og saa
var det kun hans Skygge, der gik Døren forbi: det
haster ikke; men en Dag, naar Alt er vel, og Ingen
tænker paa Noget, træder han pludselig ind og gjør sin
Pligt. Deraf have fornuftige Folk saa uddraget den Lære,
at vi, for ikke at blive overraskede, altid maae være beredte
paa Døden. Det er en god og sund og praktisk
Lære, naar vi blot vidste, hvori Beredelsen egenlig skal
bestaae.

At berede sig paa Døden bliver i borgerlig Praxis
almindeligvis forstaaet om, hvad man kalder »at beskikke
sit Hus», være sig ved at samle en Formue og gjøre sit
Testament eller ved at sætte i Enkekassen eller i Livrenten,
ved, kort og godt, at sikkre de Efterlevende
et Udkomme. Naar dette er vel besørget, saa siger en
borgerlig Bevidsthed trøstig ved sig selv: »nu er jeg beredt,
nu kan jeg døe hvad Dag det skal være, thi Regnskabet
er i Orden«. At en saadan Beredelse kan være
i høi Grad agtværdig og nødvendig, maae vi indrømme;
det røber forstandig Omhu og et kjærligt Sind at sørge
for sine Nærmeste; en Mand, der har opfyldt disse Betingelser,
kan gaae hæderlig ud af Verden og faae en
hæderlig Begravelse: Ligtalen vil ikke fortie hans Ros.

Men dersom det, at berede sig paa Døden, udeluk%
% --- p. 195 ---
\setcounter{footnote}{0}%
\opage{195}kende bestaaer i at faae de Efterlevende forsørgede, saa
maa der jo være Mange, der ere fritagne for at berede sig;
de Enlige behøve det ikke, thi de have Ingen i Verden
at forsørge; de Rige heller ikke, thi Arvingerne ville nok
vide at forsørge sig selv; de Fattige, ja de Fattige kunde
haardt trænge til en Beredelse, men de have desværre
Intet, hvormed de kunne berede sig. At berede sig paa
Døden maa, naar Fordringen skal gjælde for Alle, nødvendig
forstaaes om at berede sig paa en Tilstand efter
Døden. Men --- efter Døden? denne Tilstand er jo aldeles
ubekjendt; hvorledes kan man forberede sig paa det
aldeles Ubekjendte? Man kan forberede sig paa noget
Virkeligt, noget forud Bestemt, Noget, hvorom man veed,
at det nødvendig vil indtræde; man kan f. Ex. være
fattet og forberedt paa et Jordskælv, en Solformørkelse,
en uundgaaelig Krig, Tabet af sin Formue o. s. v. Jo
større Visheden er om det Kommende, det være sørgeligt
eller glædeligt, desto mere ligefrem er Forberedelsens
Opgave; men omvendt, jo mere ubestemt og ubestemmeligt
det Tilkommende er, desto pinligere og misligere er
Forberedelsen. At forberede sig paa det Mulige er
allerede langt vanskeligere end at forberede sig paa det
Nødvendige; men at forberede sig paa alle mulige Muligheder,
og det aldeles Ubekjendte seer jo ud som alle mulige
Muligheder, er umuligt. Det er vel især den sidste Betragtning,
der skal tjene til Undskyldning for de Mange,
der, opfyldte af Livets Travlhed, hverken have Tid eller
Lyst til at gruble over hint gaadefulde Ubekjendte ---
efter Døden.

Og hvad siger saa Døden? Ja, Døden siger naturligvis
Ingenting; og dog, naar den Huløiede nærmer sig
% --- p. 196 ---
\setcounter{footnote}{0}%
\opage{196}og stirrer paa En, er det livagtig som om Tausheden
selv blev talende, som om den virkeligt sagde Noget.
»Efter Døden! Menneske, du lyver for dig selv, naar du
bilder dig ind, at du slet Intet veed og Intet kan vide
om Tilstanden efter Døden. Du kalder mit Væsen gaadefuldt
og gaadefuldt det Land, hvorhen Reisen gjælder;
just derfor raader jeg dig, før du gaaer ombord paa mit
Udvandrerskib, at gjøre mit nærmere Bekjendtskab. Giv
mig et Navn, kald mig, hvad du vil, en raa Naturmagt,
en Sjælenes Skarpretter, en Aand fra Helvede, et Sendebud
fra Himlen; men Eet fordres ubetinget: du skal sige,
hvad du antager mig for; du skal have en Dom, en afgjørende Dom
om mig, du skal; thi din Dom om mig er
din egen Dom om dig selv: grav i dit Inderste, saa finder
du Nøglen til min Gaade«.

At Menneskets Dom om Døden er en Dom om hans
eget Selv, kan ved lidt Eftertanke let gjøres tydeligt.
Den, der siger, at Døden slukker Bevidsthedens Lys og
tilintetgjør Livet, han siger jo med det Samme, at hans
eget Selv, det Væsen, hvorpaa hans Bevidsthed, hans
Følelser, Forestillinger, Erindringer o. s. v. beroe, forsvinder
i Døden. Det er Selvet, der seer i Forkrænkelighedens
Flod sit eget Billede og sin egen Lov; det er
Selvet, der styrter sig i Floden; det er Selvet, der dømmer
sig selv til Døden. »Men den videnskabelige Fysiologi
har jo udtrykkelig bevist, at Selvhedsloven kun er
en organisk Naturlov«. Falsk, Falsk! Fysiologien veed
Intet, slet Intet om Selvhedsloven. De Love, hvortil
Fysiologien henfører alle Virksomheder af Liv og Sjæl
i det menneskelige Legeme, ere lutter fysiske Love,
Love for det Selvløse. Det Eneste, Fysiologien kan
% --- p. 197 ---
\setcounter{footnote}{0}%
\opage{197}gjøre for den, der vil æske dens Dom i Udødelighedsspørgsmaalet,
er at bestyrke ham i den Indbildning, at
Sjælen, der ikke kan tænkes uden Legeme, umulig kan
danne sig et Legeme af høiere Orden, naar det jordiske
er opløst. Men for at bestyrkes i denne Indbildning,
behøver man ikke at studere Fysiologi; man slipper nemmere
til Resultatet ved ganske ligefrem at give sig Sandsernes
Vidnesbyrd i Vold. Homer er ingen videnskabelig
Fysiolog, men han har, hvad Enhver med aaben Natursands
maa have, et stærkt, overvældende Indtryk af, at
Legemet er Sjælens Virkelighed. Sjælen forsvinder som
»en Skygge, en Drøm«, naar den forlader »de hvidlige
Knogler». Betegnende er det Overraskelsens Udbrud,
hvormed Achill's Skygge ved Indgangen til Underverdenen
tiltaler Odysses: »Ætling af Zeus, Laértes’s Søn, Odysses
du snilde! Hvi har du dristet til Hades dig ned, hvor
de Døde har hjemme, Væsener uden Bevidsthed, udlevede
Menneskers Gjenfærd?« At Mythen giver Skyggerne en
Tilværelse i Hades, er en dunkel Antydning af Selvhedslovens
Gyldighed. I samme Øieblik som den af Sandseindtryk
blændede Forstand vil gjøre Sjælen til et Intet,
reagerer den umiddelbare Selvfornemmelse; Fantasien gjør
igjen Forstandens Intet til Noget og holder det flygtige
Noget fast som en Skygge.

Den græske Forestilling om Skyggeriget er langt fra
at være forsvundet med Troen paa den græske Mythologie;
den gamle Mythe lever endnu i tusinde og tusinde Sjæle,
ængster dem i Smug og fylder dem med hemmelig Gru
for Døden; thi Mythen er i Sindets Dunkelhed, i det uklare
Selv, der er fremmed for sit eget Væsen. Hvad Under,
at et Menneske midt i Sandseverdenens Herlighed har en
kvælende Forudfølelse af Skyggeriget, saalænge hans
% --- p. 198 ---
\setcounter{footnote}{0}%
\opage{198}egenlige Selv endnu ikke er synderlig mere end en Skygge.
At leve i en saadan Mellemtilstand er utaaleligt; det er
næsten, som at være levende begravet.

Der er da ogsaa hos de Fleste en Bestræbelse efter
at befrie sig fra Mellemtilstanden; men Midlerne, der
anvendes, ville ikke altid svare til Øiemedet. Der er
dem, der ansee Adspredelser og mange Slags udadgaaende
Interesser for et godt bedøvende Middel, og der er dem,
hvem det daglige Arbeide i Forening med Næringssorger
og alskens timelige Bekymringer giver fuldt op at tænke
paa; men alt dette er usikkert. Et uventet Tilfælde, et
Dødsfald, der ret gaaer En til Hjerte, kan pludselig formørke
Synet og mane Skyggerne fra Hades frem i den opskræmmede
Fantasi. Ligesom i de varme Lande en i Dyndet nedgravet
Slange eller Krokodil kan dukke op og forskrække
den Veifarende, saaledes kan Spørgsmaalet, det uhyggelige
Spørgsmaal om Tilstanden efter Døden, uventet reise sig
og sætte Sind og Tanker i Røre. Skal Frygten overvindes,
maa Sagen tages fra Grunden af; denne Tanke
synes da ogsaa med den fremadskridende Oplysning at
vinde mere og mere Indgang.

Den moderne vidende Selvbevidsthed er allerede saa
temmelig paa det Rene med sig selv om, at dette med
Sjælens Udødelighed eller, som vi i denne Sammenhæng
helst udtrykke det, med Selvets Evighed, er en gammel
Fabel, et Børneeventyr, der hører hjemme i Ammestuen.
Selvet, et organisk Naturprodukt, opstaaer i Tiden og
forgaaer i Tiden, dette er en af den metafysiske Selverkjendelses
store Opdagelser, Udbyttet af Aarhundreders
Grublen, et frugtbringende Resultat for klare, ja for
vandklare Hoveder. Er heri noget Bedrag, da hidrører
% --- p. 199 ---
\setcounter{footnote}{0}%
\opage{199}det ikke fra Videnskaben --- Videnskaben er ganske
uskyldig --- men fra Selvet, der har mistydet og misbrugt
Videnskaben. Det er ikke en fysisk Naturlov, men en
Selvhedslov, der ligger til Grund for Dommen. Det er
Selvet, der i denne Sag udelukkende har med sig selv at
gjøre; det er Selvet, som for at komme ud af den kvælende
Uvished, med fuld Fatning bereder sig til Døden,
gjør sit Testament og testamenterer sig selv Tilintetgjørelse.

Der er vel dem, der ansee Valget af en saadan Udvei
for et Tegn paa Fortvivlelse, da Tanken om en Tilintetgjørelse
er saa gruelig. Dette er en temmelig naiv Betragtning;
thi den, der virkelig er forvisset om sin Tilintetgjørelse,
har Intet at grue for. Der tales om Dødens
Nat, om Gravens Mørke, om Ensomheden i »det kolde,
snevre Kammer« o. s. v., lutter barnagtige Fantasier.
Der er intet Mørke i Graven, ingen Ensomhed i Kisten;
for den i Døden Tilintetgjorte er der, bogstavelig talt,
hverken Lys eller Mørke, hverken Selskab eller Ensomhed:
den Døde selv er ikke engang i Graven, thi han
er slet ikke.

Der er dem, der mene, at hver den, der har dømt
sin Sjæl til Tilintetgjørelse, har begaaet et aandeligt
Selvmord, en Forbrydelse, der tilsidst maa medføre en
saadan Livslede, at Forbryderen ender med at tage sig
selv af Dage. Ogsaa i denne Synsmaade skjuler sig en
ikke ringe Misforstaaelse. Vistnok er Selvmord altid en
mulig Udvei for den, der, død for Evigheden, i sit Inderste
lever paa et Intet; men det forhøier igjen Selvfølelsen at
vide, at Mennesket ubetinget har sit eget Liv i sin egen
Haand. Det er saa langt fra, at Tilintetgjørelsestanken
% --- p. 200 ---
\setcounter{footnote}{0}%
\opage{200}uden videre skulde give Individerne Ulyst til at leve, at
det Modsatte langt snarere er Tilfældet. Den, der veed,
eller dog mener at vide, med sig selv, at Alt er ude med
Døden, kan ret faae Lyst, umaadelig Lyst til at bruge
og nyde dette korte nærværende Liv. Har Tilintetgjørelsestanken
først fra Grunden af skilt Selvet ved alt det
Evige og gjort det til et timeligt Væsen, saa maa jo
dette timelige Væsen sænke sig dybt i det Timelige, faae
en graadig Appetit paa det Timelige, nære sig deraf og
leve deri som i sit eneste rette, oprindelige Element.
Kun den forstaaer at leve, som ret forstaaer at tilfredsstille
Livets Fordringer. Den høieste Tilfredsstillelse er
timelig Selvtilfredsstillelse; den timelige Selvtilfredsstillelse
er Nydelse. For at indvinde Nydelsens Betingelser, maa
man arbeide. Arbeidet er Midlet, Nydelsen Maalet; Arbeidet
er Loven, Nydelsen Evangeliet; det er Tidens Religion.
Det Nærværende er Alt, efter Døden er Intet;
det er Tidens Filosofi. Nutidens Religion gaaer fuldstændig
op i dens Filosofi: dens Tro er i Sandhed Eet
med dens Viden.

Men Ideerne, ja dem havde vi nær glemt; Ideerne
maae naturligvis med, naar det Hele skal være komplet.
De store Verdensideer i Politik og Historie, Skjønhedsideerne
i Poesi og Konst, det Sandes Ideer i Videnskaberne,
alt dette hører væsenlig med i det nærværende
virkelige Liv. Ideerne have det umiskjendelige Fortrin,
at de paa een Gang ere baade jordiske og overjordiske;
de afgive en Garanti for, at det overalt nærværende
Evige i alle Maader er beredt til at tjene det Timelige.
Den moderne Evighed er liberal; det er en stor Behagelighed
ved de evige Ideer, at de i Grunden ere saa beskedne.
% --- p. 201 ---
\setcounter{footnote}{0}%
\opage{201}Timelighedens Kreditorer ere ikke beskedne; tvertimod,
det er paatrængende, plumpe Realister, som man ikke
kan blive kvit, før man har betalt til sidste Skilling.
Men de evige Ideer melde sig ikke som Kreditorer; det
er humane, uegennyttige Væsner. Forstaaer man at bruge
dem, godt! saa lade de sig bruge, af Mange endog med
ubegribeligt Held og forbausende, glimrende Virkning. Har
man ingen Brug for dem og ønsker ikke videre at indlade
sig med dem, godt! saa trække de sig tilbage med et Smil,
og man mærker ingen Bebreidelse. Det skal da ikke
kunne siges, at det Selv, der vil tilintetgjøres i Døden,
har udelukket sig fra Evigheden, langtfra! Medens andre
Sjæle maae nøies med Haabet, det Haab, at de efter
Døden omsider skulle indgaae i Evigheden; har dette Selv
gjort Haabet overflødigt, og er kommen Døden i Forkjøbet,
har, istedenfor selv at indgaae i Evigheden, ladet Evigheden
komme sig imøde ved at indgaae i Tiden.

Vi have, siges der, fuld Erstatning for Udødeligheden,
naar vi ville leve for Ideer. Men for at fuldkomme Selvhedsloven,
maae vi vel lægge Mærke til, hvad det egenlig vil
sige: at leve for Ideer. Naive Sjæle, der let lade Begeistringen
løbe af med sig, opfatte gjerne Sagen, som om
det, at leve for en Ide, var det Samme som at gjøre sig
til Ideens Redskab, offre Alt for Ideens Skyld, have sit
Liv i Ideens Tjeneste. Det er en skjøn Talemaade, men
en skjæv Opfattelse af Forholdet. Det kan, paa vort
valgte Standpunkt, umulig være Ideen, der ubetinget skal
herske over Selvet; det maa ifølge Selvhedsloven nødvendig
være Selvet, der skal beherske Ideen. Lad os begynde
med Kjærligheden; de Begeistrede sige, at de elske
Ideen for dens egen Skyld; det lader sig høre, forsaavidt
% --- p. 202 ---
\setcounter{footnote}{0}%
\opage{202}de dermed mene, at det er Ideen, der forskaffer dem de
største Nydelser, giver deres Liv den rigeste Solglands,
og under alle Omvexlinger af Kamp og Hvile, Sorg og
Glæde yder dem den høieste Selvtilfredsstillelse. Man
kan da elske Ideen for dens egen Skyld, som man kan
elske en ung Pige for hendes egen Skyld: Kjærligheden
varer, saalænge Selvtilfredsstillelsen varer; Selvtilfredsstillelsen
varer, saalænge Lysten varer, og Lysten ---
saalænge Nydelsen varer. Men heller ikke et Øieblik
længer. Tag Nydelsen bort, saa gaaer ogsaa Lysten bort,
og naar Lysten er borte, hvad skal saa give Kjærligheden
Næring? Pligten! Det er et Selvbedrag. Er der intet
Evighedsanlæg, ingen Spire til evigt Liv, i Selvet, hvor
skulde det da kunne føle sig forpligtet til at opoffre sig
for Noget, som er dets Natur imod? Det er en uhyre
Modsigelse at ville forpligte et timeligt Selv til at arbeide
for evige Formaal. Retfærdighedens Ide f. Ex. kan
intet Menneske med klar Forstand tillægge ubetinget
Gyldighed; Retfærdigheden staaer i Klogskabens Tjeneste
og bliver at anvende efter Omstændighederne. Retfærdighedspligten
er en Tvangspligt, og enhver Tvangspligt ophører
med Tvangen. Den Svage maa være retfærdig
mod den Stærke, men ikke omvendt. Den mindre Stat
bør til Punkt og Prik opfylde sine Forpligtelser og vise
sig retfærdig lige over for den større; men behøver en
Stormagt derfor altid at være retfærdig lige over for en
lille Magt? Det Samme gjælder om Godhedens Ide. Enhver,
der gjør det Gode, maa, hvis han er ærlig, tilstaae,
at han egenlig ikke gjør det for det Godes Skyld, men
for sin egen Skyld, fordi det yder ham Tilfredsstillelse.
Hvis Nogen derimod ubetinget vilde gjøre det Gode for
% --- p. 203 ---
\setcounter{footnote}{0}%
\opage{203}det Godes Skyld, da maatte han jo ogsaa være villig til
at gjøre det i det Tilfælde, at det for ham selv var et
Onde, var ham modbydeligt; men det er en skrigende
Modsigelse.

Der gives en filosofisk Ethik, som, ihvorvel den
ganske oprigtig fornægter Selvets Evighed, Sjælens Udødelighed,
ikke desmindre i metafysisk Præketone taler med
Salvelse om det Godes og Sandes ubetingede Krav i
Samvittigheden, om Menneskets Forpligtelse til at gjøre
det Gode endog indtil den yderste Grad af Selvfornægtelse
o. s. v. Det er en tvetydig Ethik, der taler med
to Tunger. Er Evighedsspiren indtil sidste Trevl rykket
ud af Selvet, da kan der umulig være nogen Evighedsdom
i Samvittigheden. Et timeligt Væsen kan lige
saa lidt forpligtes til at sætte Livet ind paa et evigt
Gode, som en Fisk kan forpligtes til at synge eller en
Abe til at studere Generalbas. Ud fra det Timelige gaaer
Idebevægelsen uophørlig gjennem det Evige tilbage i det
Timelige, denne Kredsbevægelse er Selvhedslovens exakte
Formel: det Timelige er Udgangspunktet, og det Timelige
er Maalet; Evigheden, det flygtige Mellemværende,
faaer alene Betydning ved at kaste Glands paa det Timelige
og afgive Vækkelsesmidler for det Nærværende.

Bilder man sig ind at kunne kyse Livet af denne
Anskuelse ved at hyle over, at det Hellige trædes under
Fødder, at de sædelige Magter haanes, og Grundvolden
for det menneskelige Samfund undergraves, da vil en
korybantisk Larmen mod ethisk Maaneformørkelse være
uden al Betydning. Her er ikke Grund til at larme,
men alvorlig Opfordring til at tænke. Det følger ingenlunde af den opstillede Synsmaade, at Lovlighed, Orden
% --- p. 204 ---
\setcounter{footnote}{0}%
\opage{204}og Sædelighed skulde gaae under eller endog blot lide
noget virkeligt Afbræk. Endskjøndt det menneskelige
Selv ikke er evigt, kan det dog, saalænge det lever, ikke
ophøre at være fornuftigt. Det Fornuftige er, at Mennesket
i Bevidstheden om sig selv har en Bevidsthed baade
om det Timelige og om det Evige, om det Timelige som
Realitet, om det Evige som Skin. Sandheden er, at det
Evige, som, naar det tages for sig og tillægges Selvgyldighed,
er en tom Tanke, en Fantasi, ikke desmindre er
af overordenlig Virkning, naar det bøies ind under det
Timelige og gaaer i Eet med det Endelige. Det gaaer
med det Evige i Bevidstheden som med Rummet i den
ydre Natur. Rummet i og for sig, det rene uendelige
Rum, er ingen Ting, er kun den rene, uendelige Tomhed;
og dog kunne de virkelige Legemer ikke være til
uden Rummet. Saaledes med det Evige. Evigheden i
sig selv er ingen Ting, og dog kan det fornuftige Menneskeliv
umulig være til uden Evighedstanken. Tyngdeloven
viser, hvorledes Afstandene, altsaa Rumsforholdene,
faae Indflydelse paa Kræfternes Virksomhed; Selvhedsloven
viser, hvorledes Ideerne, altsaa Evighedsforholdene, faae
Indflydelse paa Fornuftens Virksomhed. Sandhed og
Godhed, Pligt og Ret, disse for Samfundslivet uundværlige
Betingelser, ville ikke forsvinde fra Jorden, fordi den
tomme Forestilling om Sjælens Udødelighed omsider forsvinder
for Selvet.

Grundvolden for al Sædelighed er Oprigtighed, Ærlighed
og Redelighed. Den store Masse af flittige, adstadige,
skikkelige Borgere, som i alle Stater ere den
bestaaende Ordens bedste Støtter, giver et talende Vidnesbyrd
om, at Begreberne: Ret og Pligt, ere laante fra det
% --- p. 205 ---
\setcounter{footnote}{0}%
\opage{205}Timelige og kun finde Anvendelse paa det Timelige.
Hvilke høiere Anelser den Enkelte kan have om en sand
Hellighed, en fuldkommen Retfærdighed, en evig Gjengjældelse,
enten han nu har disse Anelser fra Præst og
Degn eller fra sit eget Indre, saa maa han dog, naar han
vil være oprigtig, tilstaae, at slige ideale Forestillinger
ere udelukkede fra Børsen, fra Torvet, fra det store Verdensmarked,
at de slet ikke gjælde i Handel og Vandel, at
de, rentud sagt, ikke finde nogen virkelig Anvendelse paa
virkelige Ting. »Enhver er en Tyv i sin Næring«. Dersom
Alle kunde komme overens om at gjøre denne Tilstaaelse,
vilde jo vistnok en Del høitravende Talemaader
forsvinde, men til Gjengjæld vilde der komme langt mere
Overensstemmelse mellem Ord og Gjerning, langt mere
Paalidelighed, Jævnhed og Naturlighed ind i de sociale
Forhold.

Ganske vist har det noget for den umiddelbare Følelse
Oprørende, at Døden, ogsaa i sædelig Forstand, gjør
Alle lige. Den Ædle døer, og Skurken døer, og Døden
slaaer en sort Streg over Regnskabet, og sletter Forskjellen
ud mellem Dyd og Last. Denne kjølige Upartiskhed
gjør imidlertid, naar vi ret besinde os, et storartet moralsk
Indtryk; den dæmper det Selviske og afkjøler Lidenskaben.

Det er jo dog en Urimelighed, at et Menneske skal
gjælde for uegennyttigt, som her i Verden gjør det Gode
uden Løn, og hvorfor? Fordi Personen ikke vil lade sig
nøie med den Smule Løn, det Nærværende kan byde ham,
men skruer Prisen op til en saadan Høide, at der maa
en hel Evighed til for at udbetale ham Lønnen. Hvilken
skjøn Uegennyttighed! En stakkels fattig Arbeider, som
% --- p. 206 ---
\setcounter{footnote}{0}%
\opage{206}forsmaaer en Løn paa ti Rigsdaler om Mandagen i det
Haab, at den ved en Strike vil stige til det Dobbelte,
naar han blot kan vente til næste Løverdag, er ikke saaledes
uegennyttig. Den lovpriste Uegennyttighed, der
sender de opskruede Regninger ind til Evighedens Bureau,
er en topmaalt Egennytte, en underfundig Strike af
umættelig Selviskhed. Mod denne Strike maa Døden, den
stærke Arbeidsherre, holde Stand: vi ere jo alle Dødens
Dagleiere --- fra Dødens Kontor skeer ingen Udbetaling;
efter Døden er der ingen Løn, thi efter Døden er der
Intet.

Ogsaa fra den anden Side er Tanken, ret beseet,
opløftende stor og moralsk. Naar den Uretfærdige ---
Røveren, Bedrageren, Morderen --- døer, saa føle alle
Retsindige, hvor velgjørende det er, at Døden holder Justits
og sætter Forbrydelsen en Skranke. Saa kysse de
Retsindige Dødens Haand og prise dens gode Gjerninger.
Men de overvættes Dydige, hine uegennyttige Sjæle, der
vente paa den store Løn, ere ikke fornøiede med en saa
simpel Retfærdighedshandling. Som et Tillæg til Fortjenesten,
et Extragratiale til Lønnen hisset, forlange de
ovenikjøbet den raffinerede Nydelse, at medens de selv
belønnes gjennem en hel Evighed, skal den Ugudelige i
deres Paasyn straffes gjennem en hel Evighed, skal uden
Naade eller Pardon gjennemlide de evige Helvedkvaler,
hvilken umættelig Hævnfølelse! Det er dog et skjønt og
ophøiet Træk i Dødens Karakter, at den forsmaaer at
give en saadan Lidenskab Næring. Den kolde Død udslukker
Hævnens Ild: efter Døden er der ingen Straf, thi
efter Døden --- er der Intet.

I skarp Modsætning til Tilintetgjørelseslæren staaer
% --- p. 207 ---
\setcounter{footnote}{0}%
\opage{207}nu Evangeliet om de Dødes Opstandelse, et Aabenbaringens
Ord, der uadskillelig er Et med Evangeliet om
Kristi Opstandelse. Enten --- Eller; enten troe paa
Kristi Opstandelse og forkaste Tilintetgjørelseslæren, eller
forkaste Evangeliet om den Opstandne og troe paa Tilintetgjørelsen;
et Tredie gives ikke. Med en overfladisk
Betragtning skulde man rigtignok mene, at der ikke blot
gaves et Tredie, men et Fjerde, et Femte, et Sjette, saa
mange Anstrengelser gjøres der, for at undgaae »de
store Yderligheder« og indrette sig en Udødelighedslære,
mere læmpet efter Tidens og en oplyst Fornufts Tankegang.
Men det fører til Intet; det \emph{er umuligt}. De mange
halvvidenskabelige Udødelighedsprojekter, som den rationalistiske
Oplysning har bragt for Dagen, ere i deres
inderste Grund ikke et Haar bedre end de mange uvidenskabelige
Projekter med Perpetuum mobile, som Naturlærens
halvstuderede Røvere bringe for Dagen. Begge
Slags Projekter have det tilfælles, at de staae i Strid
med Lovene, hine med Aandens, disse med Naturens
Love.

Selvhedsloven er Aandens Grundlov. Det er denne
Lov, hvortil ethvert Spørgsmaal om Tilstanden efter Døden
maa føres tilbage; det er denne Lov, der maa veilede
os ved Bedømmelsen af de mange Lignelser, Billeder
og Symboler, hvorunder det Hinsidige bliver os beskrevet
i det guddommelige Ord.

I Tilintetgjørelseslæren er Selvhedsloven, hvad ovenfor
er viist, bestemt som en Timelighedslov; ifølge Evangeliet
bliver den, hvad vi nu skulle see, afgjørende at bestemme
som en Evighedslov; Eet af To: \emph{enten} maa Sjælens
Natur opfattes saaledes, at Selvet, udsprunget af det
% --- p. 208 ---
\setcounter{footnote}{0}%
\opage{208}Selvløse, det Materielle, efter en timelig Bevidsthed om
sig selv og det i Fornuften Almene vender tilbage til sit
Udspring i det Materielle og forsvinder i det Selvløse;
eller Sjælen maa opfattes saaledes, at den, oprunden
af det guddommelige Selv, altsaa evig efter sit Princip,
ubevidst indtræder i Timeligheden, at den paa timelige
Betingelser kan udvikle Bevidstheden om sit Væsens
Evighed, og, beriget med Udviklingens Udbytte, vende
tilbage til sit Ophav. Det vil nu vise sig, at alt det,
Evangeliet med guddommelig Myndighed meddeler os om
de Dødes Opstandelse, hverken er Mere eller Mindre, end
hvad der netop kan tjene til at vække en levende Anskuelse
og personlig Forvisning om Selvhedslovens Opfyldelse.

For at forstaae dette, er det nødvendigt at aflægge
gamle Fordomme og feie det mugne, filtrede Væv af
theologisk-dogmatiske Forestillinger ud af Bevidstheden.

Naar En er i Færd med at skulle reise til Amerika
eller til Ny-Holland, er det i sin Orden, at han paa enhver
mulig Maade søger at skaffe sig de paalideligste
Underretninger om den ham ubekjendte Verdensdel, dens
Klima, dens Floder og Bjerge, Dyr og Planter, om Indbyggerne,
hvoraf de leve, hvorledes de gaae klædte o. s.
frdl. Men paa denne Maade kan man ikke skaffe sig
Underretninger om den tilkommende Verden, om Tilstanden
efter Døden. Men, siger man, Kristus og Apostlene
maae jo vide Besked; lad os holde os til deres Ord,
saa faae vi paalidelig Underretning. Paa denne Maade
have Theologerne fra Begyndelsen af gjort deres til at
befæste Menighederne i den Fix-Ide, at den tilkommende
Verden var et Slags overjordisk Australien eller Ny-%
% --- p. 209 ---
\setcounter{footnote}{0}%
\opage{209}Holland, hvorfra man ved Hjælp af Skriftsteder og visionære
Betragtninger, at sige, naar de behandledes paa
dogmatisk videnskabelig Maade, var i Stand til at forskaffe
sig de allerpaalideligste Efterretninger. Hvorledes
det forholder sig med disse Efterretninger, kan sees paa
hvert Blad af Dogmehistorien.

Det har kostet Anstrengelser at udfinde, hvorledes
det forholder sig med »Mellemtilstanden«, men det har
kostet endnu større Anstrengelse at faae det ganske nøiagtig
fastsat, hvorledes det egenlig skal gaae til med de
Dødes Opstandelse.

»Efter de orthodoxe Kirkelæreres, Skolastikernes og
de protestantiske Theologers overensstemmende Lære bliver
i Opstandelsen det samme Legeme restitueret med de
samme Bestanddele, som det havde i det jordiske Liv,
med Hud og Haar, Indvolde, Knogler og Tænder --- hvorledes
kunde ellers Tænders Gnidsel finde Sted i Helvede?
spurgte Hieronymus. Men nu kom man strax i Forlegenhed med Haarene. Kristus havde sagt, at ikke eet Haar
paa vort Hoved skulde gaae tabt (Luc. 21, 18). Men
hvor vilde ikke Legemet blive hæsligt, udbrød Augustin,
der satte dette Sted i Forbindelse med Opstandelsen,
naar den hele Længde af det efterhaanden afskaarne
Hovedhaar og Skæg, og maaske ogsaa Neglene, igjen
skulde voxe til det i Opstandelsen. Det var en Lykke,
at Herren paa et andet Sted havde sagt, at alle vore
Haar vare talte (Matth. 10, 30); thi heri saae man et
trøsteligt Vink om, at der med Restitutionen kun var
tænkt paa Antallet, ikke paa Længden af Haarene. Men
hvad om En i dette Liv just havde havt for faa Haar?
Legemlige Mangler af enhver Art, mener Augustin, blive
% --- p. 210 ---
\setcounter{footnote}{0}%
\opage{210}rettede ved de Fromme i Opstandelsen, Overmaal af
Fedme og Magerhed udjævnet; dog saaledes, at
den forrige Skikkelses Grundtræk, der er nødvendig
for at gjøre Personen kjendelig, ikke udslettes....
Indbildningskraften havde meget ondt ved at forestille
sig, hvorledes En, der, for at bruge Luthers Ord, har
efterladt et Ben i Ungarn, en Anden, som de vilde Dyr
have revet eller slæbt ihjel, eller som Fiskene i Havet
have spist, ved Opstandelsen igjen skulde komme til
disse Lemmer\dots{}. En mere iøinefaldende Vanskelighed
laa i det ene Tilfælde, som allerede Kirkefædrene
have opstillet, at et Menneske i en Hungersnød havde
ædt af en Andens Lig, om nemlig saa det omhandlede
Stykke Kjød skulde opstaae som en Del af den, hvem
det oprindelig tilhørte, eller af den, der siden havde assimileret
det med sig ved Fordøielsen? I sidste Tilfælde
manglede den Første et bestemt begrændset Stykke af
sit Legeme, i første mistede den Anden alle de Bestanddele
han havde tilegnet sig ved at spise Stykket. Augustin
tilkjender den oprindelige Besidder det omtvistede
Stykke Kjød og lader Gud tilbagegive den Anden de
bortdunstede Kjøddele, efter at have samlet dem af
Luften«\footnote[1]{D. F. Strauss. Den kristelige Troeslære. 2det Bind; 3die Afsnit § 103.}.

Dette blot til en lille Mundsmag for dem, hvis Tænder
let løbe i Vand, hver Gang de fornemme en Duft
fra det theologiske Kjøkken. En Opvarming af de gamle
Retter finder man endnu i hveranden »kristelig Dogmatik«,
kun at Kraften er betydelig udtyndet. Hos de

% --- p. 211 ---
\setcounter{footnote}{0}%
\opage{211}Ældre føres de eskatologiske Undersøgelser med stor Anskuelighed
og naiv Dristighed, hos de Nyere med ængstelig
Varsomhed og uendelig Flauhed: det er hele Forskjellen.

Vi holde os til Selvhedsloven. Sjælen er udødelig,
thi Selvet er evigt; dette er ingen Gjenstand for Viden,
men en uafviselig Gjenstand for Tro. »Jeg er Opstandelsen
og Livet«, siger Kristus (Joh. 11, 25—26): »hvo
som troer paa mig, han skal leve, om han end døer, og
hver den, som lever og troer paa mig, han skal i Evighed
aldrig døe. Troer du dette?« --- Ja troer du dette?
Paa Besvarelsen af dette Spørgsmaal maa da hver enkelt
Sjæl holde sin hele Opmærksomhed samlet. Bliver
Spørgsmaalet misforstaaet, da maa Besvarelsen blive derefter.
Og det bliver i Sandhed misforstaaet, naar man
opfatter Kristi Ord, som om han vilde sige: »Du skal
stole blindt paa mig, thi jeg er den udvortes, historisk
personificerede Autoritet. Naar jeg siger dig, at der er
et evigt Liv, saa er der et evigt Liv, og det maa være
dig nok«. Med denne Misforstaaelse staaer Forestillingen
om Kristi Opstandelse som en udvortes, ligefrem bevislig
Kjendsgjerning i nøie Forbindelse. »Dersom vi ikke«,
siger man, »i Skriftens Bogstav havde Sort paa Hvidt
for, at Kristus virkelig er opstanden fra de Døde, saa
vidste vi Intet, vi vidste slet ikke, hvad Haab vi turde
gjøre os om vor Tilstand efter Døden; men nu, da det
af Skriften er vist, at Kristus er opstanden, lige saa vist,
af historiske Grunde lige saa bevisligt, som at Julius
Cæsar er bleven myrdet af Brutus, og Napoleon deporteret
til St. Helena, nu vide vi, ja nu vide vi, at de
Døde skulle opstaae til evigt Liv«.

% --- p. 212 ---
\setcounter{footnote}{0}%
\opage{212}Hvor meget en saadan Forklaring strider imod Selvhedsloven,
er iøinefaldende. Selvet vil sees med sit eget
Øie, det vidende Selv med Videns Øie, det troende med
Troens Øie. Forvisningen om et evigt Liv er i dybeste
Forstand en Selvforvisning; men en oprindelig Selvforvisning
kan umulig opnaaes ved at glemme sig selv for
at høre paa en Anden, see bort fra sig selv for at see
paa en Anden, opgive sit eget Selv og kaste sig i Støvet
for en Anden. Er dette Aanden i Kristi Ord? Det er
jo en aandløs Forvrængning. Kristus taler ikke om sig
selv som om en udvortes Gjenstand; han taler om sit
Væsen. Han siger ikke: jeg har fra Himlen en Underretning
om, at der efter Døden følger et Liv og en Opstandelse;
men han sigter paa Selvet og siger til Mennesket:
jeg er det, du inderlig søger, jeg er Opstandelsen
og Livet.

Troer du dette? »Nei!« svarer den moderne Kritik.
»Jeg troer det ikke, thi jeg \emph{troer} slet ikke. Jeg er, hvad
Videnskaben kan bevidne, ikke en Troende, men en Vidende.
Troen er for Almuen, den tomme Tvivlen for de Halvdannede;
men Tankens Mænd med det klare Begreb vinde
fast Overbevisning i Viden». Man gjør uden Tvivl en
saadan Udtalelse for megen Ære ved at korse sig for
dens Dristighed, som for et gyseligt Tegn paa noget
Dæmonisk; man kommer det Sande nærmere ved at tage
Erklæringen ligefrem for hvad den er, et Tegn paa
Umodenhed i Tanken, en fuldstændig Miskjendelse af
Selvhedsloven. Selvets Evighed kan, som andensteds\footnote[1]{Jvfr. Forf. Religionsfilosofi S. 28—46; 88—97; 390—93.}
er viist, umulig erkjendes i objektiv Viden, det maa
% --- p. 213 ---
\setcounter{footnote}{0}%
\opage{213}nødvendig erkjendes i Troen. Troestanken er i sin Art
ligesaa strengt overbevisende som Videnstanken.

Den i Selvhedsloven grundende Troestanke er i alle
Dele sig selv lig, har overalt det samme Særkjende. Det
svækkede Selv maa komme til Kræfter, det dunkle, omtaagede
Selv maa klare sig i Lyset af det sande Selv.
Sjælens Indblik i sit eget Væsen er betinget af dens Indblik
i Kristi Væsen. Den Inderlighedens Kraft, hvormed
den Troende fastholder Opstandelsens Haab, beroer paa den
Følelsens Styrke, den Villiens Alvor, altsaa paa den Troens,
Inderlighedens Kraft, hvormed han formaaer at gribe det
evige Liv i Kristi Opstandelse. Udødelighedstroen kan
da forkvakles paa to Maader, idet Mennesket enten ensidig
stoler paa egen Selvforvisning eller ligesaa ensidig
bygger paa objektiv Betragtning af et udvortes historisk
Fænomen.

Overfladisk seet, skulde det jo synes, som om Mennesket,
ifølge Selvhedsloven, udelukkende eller dog allernærmest
var henvist til sig selv alene. Sagen bliver da
opfattet saaledes. Den Enkelte siger til sig selv: »Saa
sandt jeg er udødelig, maa jeg ved egen Fornuft være i
Stand til at opdage Udødelighedsspiren i mit Væsens
Anlæg. Jeg seer paa Kristus; han havde urokkelig Vished
om det evige Liv, det er sandt. Han forudsagde
sin Opstandelse og gik frivillig i Døden, deri er der noget
Opløftende, noget vidunderlig Stort og Guddommeligt.
Men det kan altsammen ikke hjælpe mig, naar jeg ikke
i mig selv, ved mig selv kan opnaae den samme Vished.
Kan jeg ikke opnaae den, hvad hjælper det saa, at jeg
læser og læser fire omtvistede og meget uoverensstemmende
Beretninger om Opstandelsens Mirakel? Kan jeg opnaae
% --- p. 214 ---
\setcounter{footnote}{0}%
\opage{214}den, saa maa jeg vel sige, at Kristus forsaavidt har
været mig en Hjælper, som han ved sit Ord og sit Exempel
har været mig en Anledning til dybere Selvforstaaelse;
men denne min Selvforstaaelse staaer ganske
uafhængig af det Historiske om hans Opstandelse«.

Mon noget Menneske paa disse Vilkaar nogensinde
har grebet det evige Liv? Det er ingen blødagtig Fromhedsmystik,
vi i Modsætning hertil ville anbefale. Kristus
er ikke en selvisk ærgjerrig Hersker, der ikke vil
unde Menneskene en Smule »moralsk Selvstændighed« af
Frygt for, at Nogen maaske skulde kunne hjælpe sig selv
uden at trænge til ham. Han har saavist Sjæle nok,
der, blot for at blive Moralloven kvit, overhænge ham
med underdanig Opvartning, sleske for ham og give
hans Naade Komplimenter. Thi det er, som ovenfor
berørt, netop det andet Brud paa Selvhedsloven, at Sjælen
blot vil forholde sig passiv overfor Aabenbareren, betragte i
from Ørkesløshed hans Lidelse, Død og Opstandelse som
et helligt Skuespil, en officiel Festforestilling, der objektivt
afslører det Eviges Hemmeligheder. Det er Ide-Nydelsens,
den æsthetiske Dogmatismes og Talemaadernes Kristendom,
en ret behagelig Kristendom, der kun har den ene Feil,
at den aldrig bringer Liv og Uforkrænkelighed for Lyset.

Mod begge Tidsvildfarelser, Moralismen og den æsthetiserende
Dogmatisme, holder Selvhedsloven kraftig igjen.
Den sidste Vildfarelse vilde snart være udryddet, dersom
den ikke havde noget opiumsagtigt Duftende, noget for
en i alt Aandeligt nervesvag Tidsalder behageligt Dulmende
ved sig. Thi Enhver veed, at det er en Vildfarelse,
og Ingen vil aabenlyst kjendes ved den. Dersom
Vildfarelsen gik nøgen omkring med Døden under
% --- p. 215 ---
\setcounter{footnote}{0}%
\opage{215}Armen, saa vilde selv Børnene paa Gaden pege Fingre
og sige: fy, hvor den Mageligheds-Kristendom dog er en
styg Vildfarelse. Men Vildfarelsen er ikke saa uopdragen,
at den skulde gaae nøgen; den veed at klæde sig, og den
har Dragter nok, flere end nogen Keiserinde.

Moralismens Vildfarelse nyder større Respekt; der er
ligesom noget Storagtigt ved den. Imidlertid kan Enhver,
der vil see lidt nærmere til, let see, at den bærer tungt
paa en Selvmodsigelse. Ja dersom Moralisten kunde gjøre
sig selv til en Kristus, kunde ved egen Kraft fylde sin
Sjæl med Sandhed, Retfærdighed og Kjærlighed til Gud
og Næsten i et saadant Maal, at han med god Samvittighed
kunde sige til sig selv: »ogsaa jeg er, før Abraham
blev, Opstandelsen og Livet«, saa vilde han visselig have
Ret. Men at tage Sagen fra denne Side falder ham vist
aldrig ind; han er jo en fornuftig Mand. Saa komme
vi da med ham ind paa en anden Vei. Han indrømmer,
at hans Væsen ikke er ubetinget Et med den evige
Sandhed, den evige Retfærdighed, den evige Kjærlighed,
men at han i Sandheds og Kjærligheds Aand føler sig
forpligtet til at stræbe, stræbe og stræbe efter Retfærdighed.
Her er Grundmodsigelsen; thi hvad jeg blot stræber
efter, det har jeg jo ikke; og hvad jeg ikke har, men
dog i det Uendelige stræber efter, er netop det, hvortil
jeg uendelig trænger.

Saa trænger da Moralisten, den fornuftige Mand,
uendelig til Forvisning om Livets Evighed, fordi han,
hvad Uroen i hans ethiske Stræben viser, uendelig trænger
til evigt Liv. Paa hvilke Vilkaar naaer han Tilfredsstillelsen,
det høie Maal, hvortil han stræber? Ikke i
% --- p. 216 ---
\setcounter{footnote}{0}%
\opage{216}og ved sig selv; thi i og ved sig selv har han kun Trangen. Heller ikke ved en Fremmed, en udenfor Staaende;
thi det, der i Aand og Sandhed skal tilfredsstille ham,
maa, om det end træder ham imøde udenfra, jo dog
forkynde sig i ham indenfra som Livets hellige Grund,
det Evighedsbrød, der mætter Sjælen, fordi det er Et med
Sjælens oprindelige Væsen.

Dette er Troens Vilkaar, Selvhedslovens Opfyldelse,
at det menneskelige Selv i uendelig Selvvirksomhed skal
drage sin Næring af Kristi Selv, skal faae Liv af hans
Liv, Lys af hans Lys, d. v. s. leve af hans Opstandelse.
Forvisningen om Kristi Opstandelse som virkelig historisk
Kjendsgjerning gaaer altsaa paa det Inderligste sammen
med den personlige Tilegnelse af Kristi Væsen. Men
denne Forvisning er sakramental. Den, der ikke har
taget Sakramentet paa, at Kristus virkelig er opstanden
fra de Døde, han kan pukke paa det Historiske saameget
han vil, og rode fysiologisk i Hjerner og Nerver saa dybt
han vil: han faaer aldrig Bugt med den ulmende Tvivl,
kommer aldrig til fuld Forvisning om evigt Liv efter
Døden.

Ved at see bort fra Kjernen og holde sig til Skallerne,
d. e. ved at opfatte Kristi Taler, ikke som en
guddommelig Opklaring af Selvets Væsen, et Livets og
Opstandelsens Lys, der i Sjælenes Nat adspreder Dødens
Mørke, men som et Slags statistisk Beskrivelse af Tilstanden
i den anden Verden, topografiske Bidrag til et
Reisekaart, saa at de Døde lig ufortrødne Turister, kunne
paa en Prik finde Veie og Stier overalt; ja, hvad mere
er, at de endnu Levende paa Forhaand næsten kunne
faae alt det Seeværdige at see baade i Himlen og Helvede,
% --- p. 217 ---
\setcounter{footnote}{0}%
\opage{217}have kristelige Dogmatikere gjort det let for den opløsende
Kritik at røre det hele Indbegreb af uligeartede
Forestillinger sammen i et selvmodsigende Virvar og nedsætte
Læren om Opstandelsen i de saakaldte Dannedes
Øine. Skal Livets Ord om de Dødes Opstandelse kun
betragtes som en Lære, da har vor Herre Kristus ikke
engang den Ære at være første Opfinder af denne Lære;
man tænke paa Mythologierne og, hvad hans nærmeste
Omgivelser angaaer, særlig paa Farisæerne. Kristus har
altsaa blot forbedret Læren, føiet et Par Paragrafer til
Læren, og dersom det maaske ifølge bibeltheologiske
Granskninger, naar de med Grundighed fortsættes endnu
et Par Tusind Aar, omsider skulde vise sig, at Evangelieberetningerne
ere nogenledes troværdige --- stadfæstet
Læren ved sin egen Opstandelse.

Et iøinefaldende Vidnesbyrd om den dogmatiske
Aandløshed ligger ligefrem deri, at Kristi forskjellige
Udsagn, naar de bogstavelig tages som faktiske Efterretninger
om det Hinsidige, let kunne fortolkes saaledes, at
de komme i aabenbar Strid med hverandre. I Lignelsen
f. Ex. om den rige Mand (Luk. 16) skeer Overgangen
fra Døden til Livet pludselig saavel til salig Glæde som
til evig Kval; og en ligesaa pludselig Overgang antydes
i Tilsagnet til Røveren paa Korset: (Luk. 23, 43). Men
ifølge Joh. 5, 28 betegnes den endelige Opvaagnen til Liv
saaledes, at de Døde først vækkes op af Gravene, naar
de »høre Guds Søns Røst«. De Hensovede maae altsaa
hvile i Jordens Skjød indtil den yderste Dag; hvisaarsag
en gammel arabisk Sekt af Kristne mente, at Sjælen
døde med Legemet og først ved den endelige Opstandelse
blev kaldt tillive med Legemet. Ligesaa øiensynlig er,
% --- p. 218 ---
\setcounter{footnote}{0}%
\opage{218}ifølge Kritiken, den uklare Vaklen mellem en jødisk
Kjødelighed og en alexandrinsk Aandelighed i den hele
Beskrivelse. Kjødeligheden er, at de Saliges Tilstand
lignes ved et stort Gjæstebud. »Mange skulle komme
fra Øster og Vester og sidde tilbords med Abraham,
Isaak og Jakob i Himmeriges Rige« (Matth. 8, 11); hvorfor
jo ogsaa Menigheden troer paa Kjødets Opstandelse.
Men Aandeligheden er, at den kjødelige Forskjel mellem
Mand og Kvinde i den Grad udslettes, at de blive til
rene Engle (Matth. 22, Luk. 20, Mark. 12); en tom.
kjedsommelig Spiritualisme, der kun lidet svarer til Kjødets
Opstandelse.

Man vil heraf skjønne, hvad godt Videnskabeligt der
kommer ud af den endeløse Brydning mellem Dogmatik
og Kritik. Overalt, hvor Kristus taler om det tilkommende Liv, taler han i Lignelser; men disse Lignelser ere
Livsbilleder, i hvilke den evige Sandhed selv vinder Skikkelse.
Saa opløftende og sjælerystende, trøsterige og forfærdelige
disse Sandhedsbilleder ere, naar de i Aandens
Lys opgaae for det indre Syn, opfylde Selvet og sætte
Sindets Kræfter i Bevægelse, saa stor er igjen Forvirringen,
naar enten en selvisk ophidset Indbildningskraft
tager Ordet forfængeligt, eller en nøgtern, aandsfattig
Forstand faaer i Sinde at stille dem op som Prospekter
til udvortes, sandselig Orientering. Hver Lignelse, hvert
Billede gjør Sandheden anskuelig med saa stor Eiendommelighed
og i saa slaaende Umiddelbarhed, at her
ikke kan være Tanke om at bringe de enkelte Billeder
til objektiv systematisk Sammenhæng. Sammenhængen
er i Grundanskuelsen, ikke i den billedlige Betegnelse.
Det er Vendepunkter, store afgjørende Hovedpunkter som
% --- p. 219 ---
\setcounter{footnote}{0}%
\opage{219}hver Gang ved særlig Anledning, i særlig Sammenhæng
blive belyste.

Et Par Exempler. Kristus siger til sine Disciple
(Matth. 19, 28): »Sandelig siger jeg Eder, at I, som
have efterfulgt mig, udi Igjenfødelsen, naar Menneskens
Søn skal sidde paa sin Herligheds Throne, skulle I og
sidde paa tolv Throner og dømme de tolv Israels Stammer«.
Men til de Skriftkloge og Farisæerne, der ville
see et Tegn af ham, siger han paa et andet Sted (Matth.
12, 41—42): »De Mænd af Ninive skulle opstaae ved
Dommen imod denne Slægt, og fordømme den; thi de
omvendte sig ved Jonas’s Prædiken, og se, her er Mere
end Jonas. Dronningen fra Sønden skal opstaae ved
Dommen, mod denne Slægt, og fordømme den; thi hun
kom fra Jordens Grændser at høre Salomons Visdom; og
se, her er Mere end Salomon«. Hvo vil nu være i
Stand til at bringe disse to Udtalelser i anskuelig Forbindelse,
naar de, vel at mærke, skulle forstaaes bogstavelig
som Beskrivelser af udvortes bestemte, i en bestemt
Orden stedfindende Optrin. Hvem skal dømme først,
Apostlene eller Niniviterne eller Dronningen af Saba?
Ikke at tale om, at Kristus selv, Alverdens Dommer,
sidder taus paa sin Trone og overlader Dommen til
Andre.

En ganske anden Betydning faae disse Udtalelser
derimod, naar vi mærke os det særlige Hovedpunkt, som
i ethvert af de store Billeder afgiver Motivet. Det er,
hvad det sidste Billede angaaer, Farisæernes Vantro, der
foranlediger Kristus til et truende Varsel mod den forhærdede
Slægt, til at minde den om, hvilken forfærdelig
Dom, den gaaer imøde. Det skal i Evigheden blive
% --- p. 220 ---
\setcounter{footnote}{0}%
\opage{220}aabenbaret, hvor høit Niniviterne, der dog kun hørte en
Jonas, og Dronningen fra Sønden, der dog kun hørte en
Salomon, i Ærefrygt for det Hellige, i Sands for det Ophøiede,
staae over en Slægt, der aabenlyst trodser den i
Kjød og Blod aabenbarede Sandhed. Dommen er paa
een Gang indvortes og udvortes; indvortes, thi det er
Slægtens egen Forhærdelse, hvorved den bliver dømt;
udvortes, thi Niniviternes Exempel og den dybe Længsel
hos Sydens Dronning ere Fænomener, der afgive Maalestok
(Jvfr. Matth. 11, 21---24). Saaledes bliver Selvhedsloven
opfyldt. Dommen er i Selvet, i hvert enkelt
Menneskes og derved i Slægtens Indre; men denne Dom
i det Indre bliver i samme Grad forfærdeligere, som
Forblindelsens og Forhærdelsens dødssvangre Følger udfolde
sig i Evighedens klare Dag, og Brøden maales med
Virkelighedens Maal.

I Udtalelsen til Disciplene er Hovedpunktet et ganske
andet. Det er det bekjendte Optrin med den rige Yngling
og Kristi strenge Ord om Rigdomssygen og de Rige, der har
foranlediget Peder til det Spørgsmaal: »see vi have forladt
alle Ting og fulgt dig; hvad skulle da vi have?«
Fuld Erstatning! lyder Svaret. Men dette Svar peger
udtrykkelig hen paa »Igjenfødelsens« sidste Katastrofe, og
faaer i Billedet Liv og Farve. Fuld Erstatning! Thi
Livet, der i Døden kun aflægger det Forkrænkelige, bliver
i Evigheden rigere, fyldigere, virkeligere, end nogen
Dødelig kan forestille sig, i det Udvortes som i det Indvortes.
Fuld Erstatning for Savn og Lidelser, for Suk
og Smerte, for Spot og Haan! Thi Endelighedens Tryk og
Byrder ere væsenlige Betingelser for Forstaaelsens Inderlighed,
for Selvudviklingens Styrke, for Kjærlighedslivet i
% --- p. 221 ---
\setcounter{footnote}{0}%
\opage{221}Tro og Haab. Fuld Erstatning i Magt og Myndighed, i
Glauds og Ære! Det er ikke et Egennyttens Maal, en
udvortes Løn for udvortes Gjerninger, som her stilles i
Udsigt; det er Livet selv, det evige Liv, en Gave saa
rig, at kun de sande Efterfølgere, de Sjæle, der paa
Selvfornægtelsens Vilkaar føle Selvhedsloven som Kjærlighedslov,
have i Kjærligheden selv Betingelsen for
at modtage Gaven. Gaven er indvortes, Aandens Gave,
Liv af Kristi Liv; men det Indvortes tager Skikkelse i
det Udvortes. Naar Kristuslivet udfolder sin Magt og
aabenbarer sin Herlighed: da skulle de, men ogsaa kun
de, hvem Æren tilkommer, faae Æren.

Den forvrængede Modsætning mellem »jødisk Kjødelighed«
og »alexandrinsk Aandelighed« falder fra dette
Synspunkt bort af sig selv. Livet efter Døden er, naar
vi holde os til Aanden i Kristi Ord, ikke et Liv af legemløse
Sjæle eller naturløse Aander; her lægges afgjørende
Vægt paa en Legemlighed af høiere Orden, paa, hvad vi
kun kunne erkjende ved at erkjende en absolut Grændse
for den menneskelige Viden, en overnaturlig Natur. Det
er Forvandlingen af den forkrænkelige Natur, hvortil Kristus
henviser i sit Svar til Sadducæerne. Det er ikke
hverken den legemlige eller den sjælelige Modsætning
af Mand og Kvinde, han vil udslette; det er Sadducæernes
paa kjødelige Forestillinger byggede Indvendinger,
han afviser. Det menneskelige Selv er ligesaa lidt en
Sjæl uden Legeme som et Legeme uden Sjæl: det er i
sin Udvikling til fuldvirkelig Eiendommelighed et sjæleligt-legemligt
Væsen. Skal nu Menneskevæsenets oprindelige
Individualitet ikke udviskes, men forklares i Opstandelsen,
da kan Ligheden med Englene umulig være
% --- p. 222 ---
\setcounter{footnote}{0}%
\opage{222}at forstaae om en Afblegen af den menneskelige Individualitet,
en Udslettelse af Mandens og Kvindens eiendommelige
Skikkelse; Stedet er, hvad Sammenhængen
udtrykkelig viser, at forstaae om Renheden og Uforkrænkeligheden
i det personlige Evighedsforhold mellem Mand
og Kvinde. Englene ere hverken Mand eller Kvinde;
en Engel har ikke Kjød og Blod. Men Kristus er Mand,
ikke Kvinde; den i Opstandelsen Forklarede har Kjød
og Blod. Kjød og Blod eller, som den Opstandne selv
har betegnet det, »Kjød og Ben«, betyder i denne Sammenhæng
ikke sandselige, forkrænkelige Elementer, men
det fuldvirkelige Menneskevæsen: Menigheden troer ikke
paa en kjødelig Opstandelse, men den troer paa Kjødets
Opstandelse.

\emph{Efter Døden --- er Dommen}; dette er Alvoren i
den hele Betragtning. Men hvad det skorter paa i Nutiden,
er aandelige Kræfter til at vække Dødstankens
Alvor saaledes, at det bliver en virkelig sand og dyb
Alvor. Den tordnende Iver, hvorved religiøse Alvorsmænd,
Bodsprædikanter, som paa Grund af det kildrende Gys,
det vækker hos Opvakte at høre om det yderste Mørke,
om Helvedpølen og Svovlsøen, naar Dampen staaer over
Andagten, og Menigheden rigtig svovles, endog kunne
blive Modeprædikanter, har ikke altid Aandens Vidnesbyrd.
Besindige Mennesker tænke som saa: »ere vi ikke
Engle, saa ere vi da heller ikke Djævle. Kommer Tid,
kommer Raad; der gives en Mellemtilstand«. Man tænke
sig et Menneske, En af dem, som, følgende Tidens Anvisning
i det religiøst Humane, saa smaat har forberedt
sig paa en Tilintetgjørelse, vaagne op efter Døden, og,
om vi maae tale saa frit, stikke Hovedet ind i Evighedens
% --- p. 223 ---
\setcounter{footnote}{0}%
\opage{223}Høisal! --- ja, han vil blive forundret, men vil han derfor
strax blive knust som en Synder? Ingenlunde! Idet den
Humane seer sig omkring og faaer Øie paa Forhørsdommeren,
gjør han en vel betænkt Undskyldning og siger
ligefrem: »jeg seer, at jeg har taget feil og beder meget om
Forladelse«. En retfærdig og human Dommer kan umulig
fordømme en Sjæl, der saa beskedent beder om Forladelse.
Den katholske Kirke tager det alvorlig med Skærsilden;
men det virker kun paa en umyndig Almue, thi Læren
om Skærsilden er kemisk udvortes. De protestantiske Kirketheologer
ræsonnere op og ned over Mellemtilstanden; der
er en evig Fordømmelse, og saa er det alligevel ingen
evig Fordømmelse. Det skal være »en dogmatisk Antinomi«
--- vi have grundig studeret Antinomien i »den
kristelige Dogmatik« --- men Antinomien er en selvlavet
Spekulation; i Antinomien er der ethiserende »Aandrighed«,
men ikke religiøs Alvor. Skal det virkelig blive
Alvor, da maae vi gaae tilbage til det religiøse Princip,
til Selvhedsprincipet, og følge Selvhedsloven strengt i Et
og Alt. Er der i Sjælene Trang, væsenlig Trang til et
evigt Liv, da maa der være dybe Grunde i Selvet. Salighedens
Muligheder maae være derinde, ja! men ogsaa
Fortabelsens Muligheder; thi det Ene betinger det Andet.
At stige ned i Selvets Afgrunde og hæve de dunkle Muligheder
op til Dagens Lys er Alvor, Livets Alvor, Dødens
Alvor: efter Døden er Dommen.
\signature{R. \emph{Nielsen}.}

\end{document}
